\section{Introduction}\label{sec:introduction}

The local--global principle in modular representation theory seeks to recover
information about representations of a finite group from normalizers of its
$p$-subgroups. Alperin's weight conjecture predicts an exact numerical form
of this principle, block by block: simple modules are counted by defect-zero
ordinary characters of normalizer quotients. Alperin formulated the
conjecture at the 1986 Arcata conference; the original account appeared in
the proceedings in 1987 \cite{Alperin1987,FLZ2023}.

Let $p$ be a prime and $G$ a finite group. For an algebraically closed field
$k$ of characteristic $p$, a block is a primitive central idempotent of
$kG$. For a block $B$, let $l(B)$ be the number of simple $kG$-modules
belonging to $B$, equivalently the number of its irreducible Brauer
characters. For character calculations we use a splitting $p$-modular
system $(K,\mathcal O,k)$: $\mathcal O$ is a complete discrete valuation
ring with characteristic-zero fraction field $K$ and residue field $k$,
and $K$ splits every subgroup of $G$. A block also denotes its unique
central lift to $\mathcal O G$. In \Cref{sec:coefficients} we justify
reduction to $k=\overline{\mathbb F}_p$ and construct such a system;
the resulting counts and block assignments are unchanged over any
algebraically closed extension of this field.

For a positive integer $a$, write $a_p$ for its largest power-of-$p$ divisor. A \emph{$p$-weight} of $G$ is a pair $(Q,\phi)$, where $Q\leq G$ is a possibly trivial $p$-subgroup and
\[
 \phi\in\Irr(N_G(Q)/Q),\qquad
 \phi(1)_p=|N_G(Q)/Q|_p.
\]
Thus $\phi$ is an ordinary irreducible character of defect zero in the quotient. The block assigned to this weight is defined using the block $b$ of the inflated character in $H=N_G(Q)$.

We specify the block-assignment convention. For $H\leq G$, coefficient restriction is the linear map
\[
 s_H:Z(kG)\longrightarrow Z(kH),\qquad
 \sum_{g\in G}a_g g\longmapsto\sum_{h\in H}a_hh.
\]
If $c$ is a block of a finite group, its central character $\lambda_c$ is the residue character of the corresponding local factor of the center. In particular it is the scalar action of the center on any simple module in $c$. We say that $b^G=B$ when
\begin{equation}\label{eq:block-assignment}
 \lambda_b\circ s_H=\lambda_B.
\end{equation}
This is the central-character definition of Brauer block induction \cite[p.~338]{Carlson1971}; see also the corrigendum \cite{Carlson1973}. In the normalizer situations arising from weights, the composite is indeed a central character; we prove this directly in \Cref{sec:group}. Let $W_p(B)$ be the set of simultaneous $G$-conjugacy classes of weights assigned to $B$.

\begin{theorem}\label{thm:main}
For every prime $p$, every finite group $G$, and every $p$-block $B$ of $G$,
\[
                  l(B)=|W_p(B)|.
\]
\end{theorem}

\Cref{thm:main} resolves the numerical blockwise Alperin weight conjecture
positively, with no restriction on the defect groups, the prime, or the
finite group. The case $Q=1$ is part of the assertion. The conclusion is
an equality of cardinalities; a canonical or equivariant correspondence
is not required for this numerical statement.

\subsection*{Local--global consequences}

The numerical equality also gives the following established consequences
of the weight conjecture. Write $\IBr_p(H)$ for the set of irreducible
$p$-Brauer characters of a finite group $H$.

\begin{corollary}[Local--global and principal-block bounds]\label{cor:block-consequences}
Let $p$ be a prime, $G$ a finite group, $B$ a $p$-block with defect group
$D$, and $b$ its Brauer correspondent in $N_G(D)$. Then
\[
 l(B)\geq l(b),\qquad
 D\text{ abelian}\Longrightarrow l(B)=l(b).
\]
For $P\in\operatorname{Syl}_p(G)$, one also has
\[
 |\IBr_p(G)|\geq|\IBr_p(N_G(P))|.
\]
Let $B_0$ be the principal $p$-block of $G$, and let $k_p(G)$ count the
conjugacy classes of $p$-power-order elements, including the identity.
Then
\[
 \begin{aligned}
 k_p(G)=3&\Longrightarrow l(B_0)\geq (p-1)/2,\\
 p\mid |G|&\Longrightarrow l(B_0)\geq 2\sqrt{p-1}+1-k_p(G).
 \end{aligned}
\]
\end{corollary}

The local block bounds are Alperin's Consequences~1--2
\cite{Alperin1987}; the principal-block bounds are due to Hung,
Sambale, and Tiep \cite[Propositions~3.1--3.2]{HST2024}.
Their weight-conjecture hypotheses are supplied by \Cref{thm:main}.
The complete derivation, including the passage to the block-free
inequality, appears in \Cref{sec:consequences}.

The block-free inequality for $p=2$ was already proved by
Malle--Navarro--Tiep \cite[Theorem~B]{MNT2023}; the local inequality for
maximal-defect $2$-blocks was proved by Feng--Mart\'inez--Rossi
\cite[Theorem~D]{FMR2025}.

The sharper classification of character-count and defect data, and the
functorial formulation of the theorem, are given in
\Cref{sec:consequences}. We first explain the historical setting and
the proof of the local--global equality itself.

\subsection*{Prior work and the proof's ingredients}

Alperin's formulation extends the notion of weights from the modular
representation theory of finite groups of Lie type in defining
characteristic to arbitrary finite groups \cite{AlperinFong1990}.
Early large-family cases include Alperin and Fong's work on symmetric
groups and on finite general linear groups in odd modular characteristic
\cite{AlperinFong1990}. Kn\"orr and Robinson expressed the conjecture
as alternating sums over chains of $p$-subgroups
\cite{KnorrRobinson1989}. Navarro and Tiep reduced the non-blockwise
conjecture to conditions on finite simple groups \cite{NavarroTiep2011},
and Sp\"ath obtained a blockwise reduction \cite{Spath2013}. These
reductions organize the problem around local correspondences for simple
groups. Among the resulting advances, Feng, Li, and Zhang proved the
blockwise conjecture for finite special linear and special unitary groups
and for finite groups with abelian Sylow $3$-subgroups \cite{FLZ2023}.
Feng, Malle, and Zhang subsequently studied generic weights for finite
reductive groups under explicit hypotheses \cite{FMZ2026}.

Other recent progress includes the inductive blockwise condition for
$F_4(q)$ at odd primes not dividing $q$, proved by An, Hiss, and L\"ubeck
\cite{AnHissLubeck2024}, and Hu and Zhou's results on inertial blocks,
including an alternative proof for $2$-blocks with abelian defect
\cite{HuZhou2025}. Huang and Y\i lmaz develop chain, Galois, and
functorial reformulations \cite{HuangYilmaz2026}. A recent preprint of
Zhang claims the inductive condition for types $\mathsf B$, $\mathsf C$
and sporadic groups \cite{Zhang2026BC}; those classes do not exhaust
the scope of \Cref{thm:main}. Our proof instead passes from a direct
block-algebra calculation to covers of curves; it does not use a
classification of finite simple groups or an inductive weight condition.

The algebraic part uses classical constructions with a specific counting
purpose. The subgroup-chain argument is related to the Kn\"orr--Robinson
reformulation. Powers modulo commutators belong to K\"ulshammer's theory
of generalized Reynolds ideals \cite{HHKM2005}, and conjugation averaging
is the group-algebra Higman trace \cite{Higman1955,LorenzTokoly2011}.
The sum of the $p$-elements detects defect-zero blocks as in Tsushima's
work \cite{Tsushima1971}.
The passage from central traces to tuples is a weighted genus-one form
of Frobenius's character-counting formula; see \cite{Klug2025} for the
surface-group formulation. We prove the exact versions needed here,
including the block projections and normalization of the tuple count.
The further assertion is uniform divisibility as the number of punctures
increases, not merely a character formula for a fixed surface.

Chen's relation between Nielsen orbit counts and degrees of moduli of
elliptic covers \cite{Chen2024} informed the geometric divisibility
viewpoint. The many-marked, high-index construction and the uniform
puncture-divisibility theorem required here are developed in this paper,
not supplied by Chen's theorem.

The geometry combines several established ideas with that divisibility
problem. High-index genus-one curves are classical objects in period--index
theory; Clark and Lacy prove arbitrary-index existence over infinite fields
finitely generated over their prime fields \cite{ClarkLacy2019}. Our
construction also prescribes marked points giving independent absolute
field differentials. The genus calculation belongs to the theory of inseparable genus
change initiated by Tate \cite{Tate1952}; see Schr\"oer
\cite{Schroer2009} for a modern treatment. Its height-one determinant
identity is a curve-level form of the canonical-bundle formula for
$p$-closed foliations \cite{Ekedahl1988}. Here the calculation must keep
additional constants and inseparable residue degrees, and the large index
turns a bounded-error estimate into an exact inequality. Finally,
maximal-immediate-extension methods of Krull and Kaplansky
\cite{Krull1932,Kaplansky1942,Poonen1993} provide a suitable valued field,
while the center-valued obstruction studied by D\`ebes and Douai
\cite{DebesDouai1997} explains the arithmetic descent issue for marked
covers. We give the required constructions and descent argument directly.

\subsection*{The mechanism of the proof}

Chain inversion first reduces the weight equality to an identity between
a defect-zero character count and an alternating sum of simple-module
counts in subgroup normalizers. A calculation with powers modulo
commutators, followed by cancellation of chains, replaces those simple
counts by ordinary-character counts. We then construct a central
group-algebra element supported on $p$-singular elements, with residue
scalars zero or one chosen to remove the defect-zero contribution from
the alternating sum. Lift it by class sums to characteristic zero and restrict
it to each chain normalizer. The alternating sum of traces of the
$p^s$th powers converges to the difference between the alternating
ordinary-character sum and the defect-zero count.
Expanding the traces gives weighted tuple counts. A further chain
cancellation leaves only tuples generating subgroups with no nontrivial
normal $p$-subgroup. The complete reduction is
\Cref{prop:group-reduction}.

Here is the counting problem that remains. Fix a finite group $J$ with
$O_p(J)=1$, where $O_p(J)$ is its largest normal $p$-subgroup, and set
$n=p^s$. We count ordered tuples $(x,y,u_1,\ldots,u_n)\in J^{n+2}$
satisfying
\[
 [x,y]u_1\cdots u_n=1,\qquad
 \langle x,y,u_1,\ldots,u_n\rangle=J.
\]
Here $[x,y]=xyx^{-1}y^{-1}$, and each puncture entry $u_i$ is
$p$-singular: its order is divisible by $p$. Prescribe a multiset of
$n$ conjugacy classes for the $u_i$, and identify tuples only by
simultaneous inner conjugation in $J$, not by permutations of their
positions. For every fixed integer $h\geq1$, we must prove that all
these counts are divisible by $p^h$ for sufficiently large $s$
(\Cref{thm:tuple-divisibility}). The threshold may depend on $p,J,h$,
but not on the multiset. This uniformity makes each weighted trace sum
tend to zero even though the number of multisets grows with $s$.

The geometric construction supplies two properties on the same marked
genus-one curve $E/F_0$ in characteristic $p$. First, every divisor
degree on $E$ is divisible by $p^s$. Second, after a cyclic splitting
extension it becomes a constant elliptic curve with independently
varying marked points. More precisely, put
$\kappa=\overline{\mathbb F}_p$, choose an ordinary elliptic curve
$A_0/\kappa$, and take $F'=\kappa(A_0^n)$. The extension $F'/F_0$
is cyclic Galois of degree $n$ and splits $E$. Over $F'$, the marked points are the
generic projection points $t_i\in A_0(F')$. If $\eta$ is a nonzero
invariant differential on $A_0$, then the pullbacks $t_i^*\eta$
form a basis of $\Omega_{F'/\kappa}$. Their independence is in this
field-differential space, not in the one-dimensional space of regular
relative differentials on the elliptic curve. After a finite extension
$F/F_0$ with $[F:F_0]_p\leq p^h$, every divisor degree on $E_F$
is still divisible by $p^{s-h}$, and at most $h$ independent
differential directions are lost over the compositum $FF'$
(\Cref{prop:twist-index,cor:index-base-change,lem:label-rank}).

For a purely inseparable cover of this curve, a calculation along
successive maps of exponent one gives a genus bound with an error controlled
by $h$ and the cover degree. The index property makes this bound exact:
after multiplication by the cover degree, the difference between the
genus term and the ramification expression is a multiple of $p^{s-h}$.
The differential calculation bounds that difference below by a negative
constant independent of $s$. Once $p^{s-h}$ exceeds the constant's
magnitude, the difference must be nonnegative. This is the rounding
step in \Cref{thm:inseparable-genus}. The resulting estimate is useful
independently of the block-algebra reduction: it applies to regular
curves over imperfect fields, retains inseparable residue degrees and
extra constants, and does not require geometric reducedness.

We lift the marked genus-one curve to mixed characteristic and pass to a
spherically complete valued field $K_s$ with value group $\mathbb Q$
and residue field $F_0$. Over an algebraic closure, the generating
tuples are precisely the monodromy data of connected $J$-covers of the
punctured genus-one curve, with the $J$-action specified. Simultaneous
inner conjugacy changes only the chosen fiber point; it does not forget
the ordering of the punctures. The lift contains all $|J|$th roots of
unity. Galois therefore preserves the local conjugacy classes as it
permutes the marked points, so it acts on the finite cover set with
the prescribed multiset. This action factors through a finite Galois
group $\Gamma$; take a Sylow $p$-subgroup of $\Gamma$.

An orbit of size less than $p^h$ would give a cover defined over an
extension $L/K_s$ with $[L:K_s]_p<p^h$. The point requiring proof is
that a fixed isomorphism class actually descends with its specified
$J$-action. The obstruction is an extension with kernel $Z(J)$.
Since $O_p(J)=1$, this kernel has order prime to $p$, and an elementary
Sylow argument supplies a descent datum
(\Cref{prop:small-orbit-descent}). Put $F=\operatorname{res}(L)$.
The full degree of $L/K_s$ equals $[F:F_0]$, so this residue extension
also has small $p$-part.

To rule out such a descended cover $Y$, extend the valuation obtained
by reducing functions on the base curve to the function field of $Y$,
and take the finite $J$-orbit of one such extension. For each valuation
$w$ in this orbit, normalizing $E_F$ in its residue field gives a
regular projective curve $C_w$. Reduction of section spaces first shows
that the degrees of the maps $C_w\to E_F$ sum to $|J|$. A second
comparison imposes vanishing above selected marked points. Together
with characteristic-zero Riemann--Hurwitz and the sharp inseparable
genus estimate, these section comparisons
force
\[
                         \chi(Y)=\sum_w\chi(C_w),
\]
where each structure-sheaf Euler characteristic is computed over its
stated ground field. The equality gives enough sections to construct a
finitely presented flat model whose special fiber is the disjoint union
of the $C_w$ (\Cref{prop:specialization-model}). Passing to a
Noetherian normal base allows us to apply connectedness. There is
therefore only one $C_w$, and its valuation is stabilized by all of
$J$. The kernel of the action on its residue field, the inertia group,
is consequently normal in $J$. The same genus comparison shows that
this kernel is a nontrivial $p$-group, contradicting $O_p(J)=1$.

Every Sylow orbit thus has size at least $p^h$. Its size is a power
of $p$, hence is divisible by $p^h$. Summing the orbits proves the
uniform tuple divisibility, and \Cref{prop:group-reduction} returns
this conclusion to the block identity.

The auxiliary fields in the geometric construction are unrelated to the
coefficient system $(K,\mathcal O,k)$ used in the block count.

\Cref{fig:proof-interfaces} displays how the two geometric inputs exclude
covers over fields whose degree has small $p$-part, and how this returns to the original
block count. \Cref{sec:foundations} records the classical inputs and
coefficient choices. \Cref{sec:group} proves the algebra-to-counting
reduction before any curve is constructed. \Cref{sec:twist} supplies the
marked high-index curve, \Cref{sec:differentials} establishes the sharp
genus estimate, and \Cref{sec:valuations} constructs the valued lift and
descends small orbits. \Cref{sec:specialization} builds the finite model,
applies connectedness, and completes both the divisibility theorem and
the block identity. Finally, \Cref{sec:consequences} proves the announced
local--global bounds and presents the character-count classification and
functorial identity with their full conventions.

\begin{figure}[!ht]
\centering
\begin{tikzpicture}[
  every node/.style={font=\small,align=center},
  stage/.style={draw=black!55,rounded corners=2pt,inner sep=7pt},
  flow/.style={-{Stealth[length=5pt]},semithick}
]
\node[stage,text width=12cm] (twist)
  {Marked genus-one curve: large divisor index\\
   and independent absolute field differentials\\
   \Cref{sec:twist}};
\coordinate (branch-top) at ($(twist.south)+(0,-9mm)$);
\coordinate (branch-split) at ($(twist.south)+(0,-4mm)$);
\node[stage,text width=5.5cm,anchor=north] (genus)
  at ($(branch-top)+(-3.2cm,0)$)
  {Differential estimate and divisor-degree divisibility\\Sharp inseparable genus bound
   (\Cref{sec:differentials})};
\node[stage,text width=5.5cm,anchor=north] (lift)
  at ($(branch-top)+(3.2cm,0)$)
  {Valued lift and marked-cover descent\\Small orbit gives an extension
   with small $p$-part of its degree
   (\Cref{sec:valuations})};
\node[fit=(genus)(lift),inner sep=0pt,outer sep=0pt] (branches) {};
\node[stage,text width=12cm,below=10mm of branches] (specialization)
  {Section comparison, finite presentation, and connectedness\\
   No such marked $J$-cover over $L$ with $[L:K_s]_p<p^h$\\
   when $O_p(J)=1$
   (\Cref{sec:specialization})};
\node[stage,text width=12cm,below=7mm of specialization] (count)
  {Every Sylow orbit has size divisible by $p^h$\\
   Uniform tuple divisibility (\Cref{thm:tuple-divisibility})};
\node[stage,text width=12cm,below=7mm of count] (blocks)
  {Chain inversion and central trace limits (\Cref{sec:group})\\
   $l(B)=|W_p(B)|$};
\draw[flow] (twist.south) -- (branch-split) -| (genus.north);
\draw[flow] (twist.south) -- (branch-split) -| (lift.north);
\draw[flow] (genus.south) -- (genus.south |- specialization.north);
\draw[flow] (lift.south) -- (lift.south |- specialization.north);
\draw[flow] (specialization.south) -- (count.north);
\draw[flow] (count.south) -- (blocks.north);
\end{tikzpicture}
\caption{Mathematical dependencies, not reading order. The covers are geometrically
connected $J$-Galois covers, with specified $J$-action over their field of
definition, branched only at the marked points and with inertia order
divisible by $p$ at every marked point. The two branches provide, respectively, the
genus estimate and the arithmetic realization of a hypothetical small
Sylow orbit. For sufficiently large $s$, their combination rules out such a cover over a field whose
degree has small $p$-part; the geometric cover degree is fixed at $|J|$.
The resulting uniform
divisibility supplies the counting hypothesis of the block-algebra
reduction, which is proved before the geometric construction and applied
at the end. No weight-counting theorem is used to construct the curves.}
\label{fig:proof-interfaces}
\end{figure}


\subsection*{Conventions}

All groups are finite. The identity is a $p$-element; a $p$-singular element has order divisible by $p$. We use $[x,y]=xyx^{-1}y^{-1}$ and write $O_p(X)$ for the largest normal $p$-subgroup of $X$. All curve degrees and coherent Euler characteristics are taken over the ground field specified at that point. A regular curve over an imperfect field is not presumed smooth or geometrically reduced. For such a projective curve $C$, we set
\[
 \chi(C)=\chi(\mathcal O_C),\qquad \nu(C)=-2\chi(C).
\]
These quantities add on disjoint unions. A valuation is written additively. A field extension is \emph{immediate} if it changes neither the value group nor the residue field.
